\subsection{Resolventes del transporte y conservación de la memoria}
\label{subsec:ii-memoria-resolvente}

La historia prolongada determina las ventanas
\(I_m=\{9m+1,\ldots,9m+9\}\), para \(m\geq0\).
Con la numeración de esta sección, \(I_m=E_m\) durante el primer ciclo.
Sea \(\sigma_m\in\{-1,1\}\) la orientación conservada por esa historia;
para \(0\leq m<12\), coincide con \(\Sigma_{12}(m)\). Se define
\[
 a_m=\sigma_m\sum_{t\in I_m}\mathbf1_{\{2,4,6,8\}}(d_t),
 \qquad |a_m|\leq9.
\]
El registro y el marco móvil se prolongan por
\[
 z_{m+1}=z_m+a_m\mathbf e_{m\bmod12},\qquad y_m=S^{-m}z_m.
\]
El estado inicial \(z_0\) conserva la historia previa; vale cero para
un prefijo inicialmente vacío. En este apartado escribimos
\(R=S^{-1}\) para el transporte del marco. La actualización
\eqref{eq:eta} se mantiene en toda ventana. La orientación y los eventos
posteriores se obtienen de la historia prolongada: el retorno del marco
no presupone periodicidad de la secuencia de eventos.

\begin{proposition}[Identidad de truncación con estado terminal]
Para $N\geq1$, sean
\[
Y_N(u)=\sum_{m=0}^N y_mu^m,\qquad
H_{\eta,N}(u)=\sum_{m=0}^{N-1}\eta_mu^m.
\]
Se cumple la identidad polinómica exacta
\begin{equation}
(I-uR)Y_N(u)=y_0+uH_{\eta,N}(u)-u^{N+1}Ry_N.
\label{mem:identidad-finita}
\end{equation}
\end{proposition}
\begin{proof}
En grados $1\leq m\leq N$, el coeficiente del miembro izquierdo es
$y_m-Ry_{m-1}=\eta_{m-1}$, por \eqref{eq:eta}.
Sus coeficientes en grados cero y $N+1$ son $y_0$ y $-Ry_N$,
respectivamente. La comparación de coeficientes prueba la igualdad.
\end{proof}

El término terminal conserva la memoria transportada al cortar una historia.
La identidad permite comparar una evaluación finita con su prolongación
manteniendo explícito el estado que queda en la frontera.


\subsubsection{Correspondencia con el desarrollo común de memoria}

El cambio de numeración de las ventanas no cambia el registro. En el
primer ciclo, las convenciones de los dos desarrollos se relacionan por
\[
 E^{(I)}_{m+1}=I_m=E^{(II)}_m,\qquad 0\leq m<12.
\]
Los signos y los eventos de cada ventana coinciden después de este
cambio de índice. La convención del transporte también coincide:
\(S\mathbf e_j=\mathbf e_{j+1\bmod12}\) y \(R=S^{-1}\).
Por tanto, la actualización \eqref{eq:eta} es la misma identidad que
la ecuación~\textup{(\HMTIIRefOriginal{mem:actualizacion})} del núcleo
común, con el mismo estado previo conservado.

La proposición~\HMTIIRefOriginal{mem:prop-serie} demuestra la identidad
de la serie generadora, su convergencia para \(|u|<1\) y la cota uniforme
de cola~\textup{(\HMTIIRefOriginal{mem:cota})}. La identidad finita
\eqref{mem:identidad-finita} añade aquí el término de frontera
\(-u^{N+1}Ry_N\): permite cortar la historia sin sustituir su estado
terminal por cero. El parámetro \(u\) cuenta, en ambos casos, ventanas
de nueve transiciones.

La eliminación de componentes de memoria conserva el operador, la fuente
y el lector según la proposición~\HMTIIRefOriginal{mem:prop-schur} y
las ecuaciones~\textup{(\HMTIIRefOriginal{mem:schur-fuente})} y
\textup{(\HMTIIRefOriginal{mem:schur-lector})}. El ejemplo del ciclo
de doce ventanas~\textup{(\HMTIIRefOriginal{mem:ejemplo-ciclo})}
distingue la compresión del resolvente del resolvente de la compresión.
La composición de tres segmentos y la transitividad de la reducción
conservan sus pruebas junto a las identidades
\textup{(\HMTIIRefOriginal{mem:cociclo-tres})} y
\textup{(\HMTIIRefOriginal{mem:schur-transitivo})}.
Estas remisiones pertenecen al desarrollo íntegro incluido en este mismo
artículo. El balance \eqref{mem:balance} y el estado terminal de
\eqref{mem:identidad-finita} conservan, en su aplicación presente,
los incrementos y la frontera del registro; no determinan por sí solos
coeficientes analíticos adicionales.

